%% Chapter 16: Complexity Theory for Approximation and Parameterization
%% Status: skeleton -- learning goals and section outline fixed, prose to be written.

\chapter{Complexity Theory for Approximation and Parameterization}
\label{ch:hardness}

Both of the previous parts leave an obvious question open: which problems admit
a PTAS, and which problems are FPT? This chapter provides the negative theory.
For approximation the key notions are APX-hardness and the L-reduction; for
parameterization they are the W-hierarchy and ETH-based lower bounds.

\section*{Learning goals}
\addcontentsline{toc}{section}{Learning goals}

After this chapter you should be able to:
\begin{itemize}[leftmargin=*]
  \item define APX and approximation-preserving reductions;
  \item define an L-reduction and prove that L-reductions compose and preserve constant-factor approximability;
  \item prove APX-hardness of \lang{Max Cut on Multigraphs} by an L-reduction from \lang{Max Not-All-Equal 3-SAT};
  \item state the PCP theorem and explain how it yields inapproximability results;
  \item define the W-hierarchy and W[1]-hardness, and give W[1]-hard problems;
  \item derive parameterized lower bounds from ETH.
\end{itemize}

\section{The class APX}
\label{sec:apx}

\todo{Write this section.}

\section{L-reductions}
\label{sec:l-reductions}

\todo{Write this section.}
% Planned content: Definition, composition, and what they preserve.

\section{Max cut on multigraphs is APX-hard}
\label{sec:maxcut-apx-hard}

\todo{Write this section.}
% Planned content: An L-reduction from \lang{Max Not-All-Equal 3-SAT}.

\section{The PCP theorem and gap reductions}
\label{sec:pcp}

\todo{Write this section.}

\section{The W-hierarchy}
\label{sec:w-hierarchy}

\todo{Write this section.}
% Planned content: FPT-reductions, W[1] and W[2], and \lang{Clique} as the canonical W[1]-hard problem.

\section{Lower bounds from ETH}
\label{sec:eth-lower-bounds}

\todo{Write this section.}

\section{Exercises}
\label{sec:ex-hardness}

\todo{Write this section.}

\begin{poolquestions}
  \poolq{18}{L-reduction for Maximum Cut on Multigraphs}Covered in \cref{sec:maxcut-apx-hard}; the definitions it relies on are in \cref{sec:apx,sec:l-reductions}.
\end{poolquestions}
